\subsubsection{3D Fixed-Wall and Resolved-FSI Models}
\label{subsubsec:3d-models}

The three-dimensional tier is the natural starting point after the governing
equations because it keeps the lumen domain, velocity field $\vu(t,\vx)$,
pressure field $p(t,\vx)$, stress tensor, and boundary decomposition of the
continuum problem. In a fixed-wall hemodynamic model, the wall geometry is fixed
or prescribed while the fluid equations are solved on the resolved lumen
\cite{QuarteroniFormaggia2004CardiovascularModeling}. This tier can supply
resolved velocity, pressure, and wall traction. A WSS output additionally
requires the stress tensor, wall normal, wall location, tangent convention, and
output operator to be declared
\cite{JohnEtAl2017WallShearStressVectorForm}. These quantities are fields on
the resolved lumen and describe the vessel wall's interaction with the fluid
domain, not the 1D state variables $A$, $Q$, and $p_{\mathrm g}$.

For the idealized stenosis used here, a fixed-wall stationary-Stokes study is
the locally feasible 3D tier. The $C^\infty$ radius profile gives an analytic
wall and repeatable tetrahedral mesh families, so resolved velocity, pressure,
and stress-derived wall-traction diagnostics can be studied without changing
the wall shape in time. This is still a fixed-wall fluid solve: it does not
recover wall displacement, wall velocity, structural stress, or interface work.

Resolved FSI changes the wall descriptor rather than merely adding a new output.
A fixed-wall model belongs to the $\mathcal W_{\mathrm{rigid}}$ or
prescribed-wall reading in Definition~\ref{def:wall-fsi-closure-catalog}; a
resolved-FSI model uses $\mathcal W_{\mathrm{resolvedFSI}}$ and adds wall
displacement, wall velocity, structural stress, interface kinematics, and
fluid--structure force balance. The elastic 1D wall law
$\mathcal W_{\mathrm{elastic1D}}$ is a reduced pressure--area closure, not a
structural FSI solve
\cite{FormaggiaEtAl2007FSICoupling,GaldiEtAl2008HemodynamicalFlows}.

The 3D/FSI tier can represent eccentricity, curvature, branches, secondary
flow, separation, and recirculation more directly than a 1D area--flow model.
That role does not make it assumption-free. A valid comparison requires matched
geometry, rheology descriptor $\mathcal R$, wall closure $\mathcal W$, boundary
data, mesh and time resolution, output operators, and discrepancy metrics. A
velocity-slab diagnostic, a traction-based WSS comparison, and a pressure-ratio
comparison are different observation problems.

The 3D/FSI tier supplies the reference vocabulary for resolved velocity,
pressure, traction, and wall motion, but it does not automatically validate a
reduced model without matched comparison data. It is not the selected broad
parametric generator in this report because the setup and computational burden
are high, not because resolved 3D or FSI models are unimportant.
